%%%%%%%%%%%%%%%%%%%%%%%%%%%%%%%%%%%%%%%%%%%%%%%%%%%%%%%%%%
%%
%%  METADATA
%%  Document type: Curriculum Vitae
%%  Author: Patrick Iglesias-Zemmour
%%  Date: May 2026
%%  Version: 1.0
%%  Repository: https://github.com/p-i-z/Diffeology-Archives
%%
%%%%%%%%%%%%%%%%%%%%%%%%%%%%%%%%%%%%%%%%%%%%%%%%%%%%%%%%%%

\documentclass[11pt,reqno]{amsart}

%%%%%%%%%%%%%%%%%%%%%%%%%%%%%%%%%%%%%%%%%%%%%%%%%%%%%%%%%%
%%  MACROS
%%%%%%%%%%%%%%%%%%%%%%%%%%%%%%%%%%%%%%%%%%%%%%%%%%%%%%%%%%

\parindent 0mm
\parskip .5ex plus 2pt

\usepackage{amssymb}
\usepackage{url}
\usepackage[T1]{fontenc}
\usepackage[utf8]{inputenc}

% For book name in Chinese
\usepackage{CJKutf8}

\usepackage{microtype}
\usepackage[cal=scr,frenchmath,uppercase = upright,greeklowercase = upright,greekfamily = didot, utopia]{mathdesign}
\linespread{1.1}

\usepackage[margin=1in]{geometry}

\newlength{\datewidth}
\setlength{\datewidth}{4cm}

\newcommand{\datedsection}[2]{\noindent\textbf{#1}\hfill\textit{#2}\vspace{.5ex}\par}
\newcommand{\dateditem}[2]{\noindent\hangindent=0.5cm #1 \hfill \parbox[t]{\datewidth}{\raggedleft\itshape #2}\vspace{.2ex}\par}
\newcommand{\pubitem}[3]{\noindent\hangindent=0.5cm #1, \emph{#2}, #3\vspace{.3ex}\par}
\newcommand{\bookitem}[4]{\noindent\hangindent=0.5cm #1 (#2). \emph{#3}. #4\vspace{.5ex}\par}

\begin{document}
  
  \begin{center}
    {\LARGE\textbf{Patrick Iglesias-Zemmour}} \\ [1ex]
    Former Researcher, CNRS (France) — 40-year career \\
    Visiting Professor, The Hebrew University of Jerusalem, Israel \\ [1ex]
    \url{http://math.huji.ac.il/~piz} \quad \texttt{piz@math.huji.ac.il} \\ [1ex]
    Born 27 December 1953, French \& Israeli citizenship
  \end{center}
  
  \medskip\hrule\medskip
  
  
  %% Research Profile
  \subsection*{Research Profile}
  
  \textbf{Founder of modern diffeology} — a geometric framework extending differential geometry to singular and infinite-dimensional spaces.
  
  Building on the axiomatic foundations laid by J.-M. Souriau (1980) and the category of \textit{espaces différentiels} introduced with P. Donato (1983), I have since developed diffeology into a mature field: fiber bundles, homotopy theory, the chain homotopy operator, the moment map (AMS Memoir 2010), orbifolds, quasifolds, the bridge to noncommutative geometry (JNCG 2018, 2021), geometric quantization by paths (AMS Memoir under review), and the definitive texts \textit{Diffeology} (AMS 2013, Beijing 2022), \textit{Lectures on Diffeology} (2025), and the forthcoming three-volume \textit{The Geometry of Motion}.
  
  My work establishes systematic bridges between diffeology and:
  
  \begin{itemize}
    \item \textbf{Noncommutative Geometry} — Two papers in the \emph{Journal of Noncommutative Geometry} (2018, 2021) show that C*-algebras of singular spaces arise naturally from diffeological groupoids.
    \item \textbf{Geometric Quantization} — Forthcoming AMS memoir \emph{Geometric Quantization by Paths} constructs the \emph{Prequantum Groupoid}, unifying Feynman's path integral with the Dirac-Kostant-Souriau program.
    \item \textbf{Homotopy Theory} — Sustained collaboration with the Japanese school (Iwase, Kuribayashi, Shimakawa) through five editions of the \emph{Building Differential Homotopy Theory} conference series.
  \end{itemize}
  
  \textbf{Unifying theme:} Universal models for singular symplectic spaces—from infinite quasi-projective spaces (Proc. AMS, 2016) to toric quasifolds (with E. Prato, JNCG, 2021) to geometric quantization by paths: the prequantum groupoid, and noncommutative geometry as the "Threshold of Analysis" (AMS memoir, under review).
  
  %% Professional Positions
  \subsection*{Professional Positions} \null\
  
  \dateditem{Visiting Professor at Università degli Studi di Firenze Sept. 2026.}
  
  \dateditem{Visiting Professor, The Hebrew University of Jerusalem, Israel}{2021–present}
  
  \dateditem{Researcher, Institut de Mathématiques de Marseille, CNRS, France}{2011–2020}
  
  \dateditem{Visiting Professor, The Hebrew University of Jerusalem, Israel}{2004–2011}
  
  \dateditem{Researcher, LATP, CNRS Marseille, France}{1998–2004}
  
  \dateditem{Researcher, UMPA, École Normale Supérieure de Lyon, France}{1992–1998}
  
  \dateditem{Researcher, Centre de Physique Théorique, CNRS Luminy, Marseille, France}{1980–1992}
  
  %% Education
  \subsection*{Education} \null\
  
  \dateditem{Doctor of Science (Doctorat d'État), 1985}{1985}
  Université de Provence. Dissertation: \emph{Fibrations difféologiques et homotopie}. Advisor: J.-M. Souriau. Board: J. Dixmier (President), D. Bennequin, J. Breuneval, P. Donato, J. Pradines, J.-M. Souriau, R. Stora. \emph{With distinction.}
  
  \dateditem{PhD (Doctorat de 3ème cycle), 1979}{1979}
  Université de Provence. Dissertation: \emph{Thermodynamique géométrique appliquée aux configurations tournantes en astrophysique}. Advisor: J.-M. Souriau. \emph{With distinction.}
  
  \dateditem{DEA Mathématiques Pures (Master), 1978}{1978}
  Université de Provence. \emph{With merit.}
  
  \dateditem{DEA Physique Théorique (Master), 1977}{1977}
  Université de Provence. \emph{Valedictorian, with distinction.}
  
  %% Recent Teaching Experience
  \subsection*{Recent Teaching Experience} \null\
  
  \dateditem{Invited Lecture Series}{Oct 2020 – Feb 2021}
  \emph{Diffeology: From Basics to Recent Developments} (17 weeks, 34 hours), Shantou University, China (via Zoom).
  Program: "Famous Teachers Overseas" (\begin{CJK}{UTF8}{gbsn}海外名师\end{CJK}).
  Materials formed the basis for the first chapters of \textit{Lectures on Diffeology} (Beijing WPC, 2025)
  
  \dateditem{Post-Graduate Course on Diffeology}{Apr 13 – May 3, 2026}
  Tel-Aviv University, Israel.
  
  \dateditem{Summer School, Turkey}{2018}
  \emph{Two-week lecture series on diffeology}, Nesin Mathematics Village.
  
  \subsection*{Books and Monographs} \null\
  
  \bookitem{1}{2026}{The Geometry of Motion (3 volumes)}{Beijing World Publishing Corp., under contract. Preprint: \url{http://math.huji.ac.il/~piz/documents/TGOM.pdf}}
  
  \bookitem{2}{2025}{Lectures on Diffeology}{Beijing World Publishing Corp., 378 pages. ISBN 978-7-5232-1841-9}
  
  \bookitem{3}{2022}{Diffeology (Revised Second Edition)}{Beijing World Publishing Corp., 529 pages. ISBN 978-7-5192-9608-7}
  
  \bookitem{4}{2018}{La Géométrie des Mouvements}{Amazon Desktop Kindle Publishing, 71 pages. ISBN 978-1-9810-3206-8}
  
  \bookitem{5}{2013}{Diffeology}{American Mathematical Society, Mathematical Surveys and Monographs, vol. 185, 450 pages. ISBN 978-0-8218-9131-5}
  
  \bookitem{6}{2010}{The Moment Maps in Diffeology}{American Mathematical Society, Memoirs, no. 972, vol. 207, 80 pages. ISBN 978-0-8218-4709-1}
  
  \bookitem{7}{2000}{Symétries et Moment}{Hermann, Paris, 200 pages. ISBN 978-2-7056-6419-0}
  
  %% Works in Progress
  \subsection*{Works in Progress} \null\
  
  \pubitem{2026}{Geometric Quantization by Paths}{AMS Memoir, submitted. Based on three preprints:}
  \url{http://math.huji.ac.il/~piz/documents/GQBP-I.pdf} (Part I: Simply Connected, 2025) \\
  \url{http://math.huji.ac.il/~piz/documents/GQBP-II.pdf} (Part II: General Case, 2025) \\
  \url{http://math.huji.ac.il/~piz/documents/GQBP-III.pdf} (Part III: Metaplectic Anomaly, 2026) \\
  arXiv:2508.11337, 2512.24627, 2601.23259
  
  %% Selected Journal Publications
  \subsection*{Selected Journal Publications} \null\
  
  \pubitem{2026}{The Boman Paradox: When Topology Plus Algebra Do Not Define Geometry}{Math. Intelligencer (2026). DOI: 10.1007/s00283-025-10500-3}
  
  \pubitem{2025}{Why Diffeology}{Arnold Mathematical Journal, Vol. 12, Issue 2 (DOI not yet active)}
  
  \pubitem{2025}{Groupoids in Diffeology}{Electronic Research Archive, Vol. 33, No. 12, pp. 7957–7973. DOI: 10.3934/era.2025350}
  
  \pubitem{2025}{On the Diffeology of Orbit Spaces (with S. Gürer)}{Differential Geometry and its Applications Volume 103, June 2026. DOI: 10.1016/j.difgeo.2026.102391}
  
  \pubitem{2024}{Editorial note to: On the motion of spinning particles in general relativity by Jean-Marie Souriau (with T. Damour)}{General Relativity and Gravitation, 56(10):127}
  
  \pubitem{2023}{Orbifolds as Stratified Diffeologies (with S. Gürer)}{Differential Geometry and its Applications, vol. 86}
  
  \pubitem{2023}{Čech–De Rham Bicomplex in Diffeology}{Israel Journal of Mathematics}
  
  \pubitem{2021}{Quasifolds, Diffeology and Noncommutative Geometry (with E. Prato)}{Journal of Noncommutative Geometry, vol. 15, no. 2, pp. 735–759}
  
  \pubitem{2021}{Every Symplectic Manifold is a (Linear) Coadjoint Orbit (with P. Donato)}{Bulletin of the Canadian Mathematical Society, pp. 1–16}
  
  \pubitem{2019}{Refraction and Reflexion According to the Principle of General Covariance}{Journal of Geometry and Physics, vol. 142}
  
  \pubitem{2018}{Noncommutative Geometry and Diffeology: The Case of Orbifolds (with J.-P. Laffineur)}{Journal of Noncommutative Geometry, vol. 12, no. 4, pp. 1551–1572}
  
  \pubitem{2016}{Example of Singular Reduction in Symplectic Diffeology}{Proceedings of the American Mathematical Society, vol. 144, pp. 1309–1324}
  
  \pubitem{2015}{Primary Spaces, Mackey's Obstruction, and the Generalized Barycentric Decomposition (with F. Ziegler)}{Journal of Symplectic Geometry, vol. 13, no. 1, pp. 55–80}
  
  \pubitem{2013}{Variation of Integrals in Diffeology}{Canadian Journal of Mathematics, vol. 65, no. 6, pp. 1255–1286}
  
  \pubitem{2011}{Smooth Lie Group Actions are Diffeological Subgroups (with Y. Karshon)}{Proceedings of the American Mathematical Society, vol. 140, pp. 731–739}
  
  \pubitem{2010}{Orbifolds as Diffeologies (with Y. Karshon and M. Zadka)}{Transactions of the American Mathematical Society, vol. 362, no. 6, pp. 2811–2831}
  
  \pubitem{2007}{Dimension in Diffeology}{Indagationes Mathematicae, vol. 18, no. 4}
  
  \pubitem{1998}{Les origines du calcul symplectique chez Lagrange}{L'Enseignement Mathématique, vol. 44, pp. 257–277}
  
  \pubitem{1995}{La trilogie du moment}{Annales de l'Institut Fourier, vol. 45, fasc. 3}
  
  \pubitem{1991}{Les SO(3)-variétés symplectiques et leur classification en dimension 4}{Bulletin de la Société Mathématique de France, vol. 119, fasc. 4}
  
  \pubitem{1985}{Exemple de groupes différentiels: flots irrationnels sur le tore (with P. Donato)}{Comptes Rendus de l'Académie des Sciences, t. 301, série I, n° 4}
  
  %% Recent Preprints
  \subsection*{Recent Preprints (Under Review)} \null\
  
  \pubitem{2026}{A Conformal Co-Symplectic Structure on the Space of Pseudo-Riemannian Geodesics}{Submitted to the Journal of Symplectic Geometry. \url{http://math.huji.ac.il/~piz/documents/ACCSSOSOPRG.pdf}}
  
  \pubitem{2025}{Does the motor cortex draw on a wire plane?}{Submitted to Biological Cybernetics. \url{http://math.huji.ac.il/~piz/documents/DTMCDOTWP.pdf}}
  
  \pubitem{2025}{Bosonic and Fermionic Singularities in Diffeology}{\url{http://math.huji.ac.il/~piz/documents/BAFSID.pdf}}
  
  \pubitem{2025}{Diffeology and Arithmetics of Irrational Tori}{\url{http://math.huji.ac.il/~piz/documents/DAAQIT.pdf}}
  
  
  %% Major International Collaborations
  \subsection*{Major International Collaborations} \null\
  
  \begin{itemize}
    \item \textbf{Japan:} Norio Iwase (Kyushu), Katsuhiko Kuribayashi (Shinshu), Kazuhisa Shimakawa (Okayama), Hiroshi Kihara (Aizu). Invited participant in all five \emph{Building Differential Homotopy Theory} conferences (2019–2026). Iwase's JSPS grant explicitly identifies me as the ``central figure'' in diffeology research.
    \item \textbf{Italy:} Elisa Prato (Firenze). Collaboration on quasifolds and noncommutative geometry since 2017; joint JNCG paper (2021).
    \item \textbf{Turkey:} Serap Gürer (Galatasaray). Collaboration since 2015; multiple joint papers on diffeology, orbifolds, and differential forms.
    \item \textbf{Israel/Canada:} Yael Karshon (Toronto). Collaboration since 2002; joint work on orbifolds and Lie group actions.
    \item \textbf{China:} Enxin Wu (Shantou). Invited lecture series (2020–21) and ongoing collaboration.
  \end{itemize}
  
  %% Invited Lectures
  \subsection*{Invited Lectures (Last 5 Years)} \null\
  
  \dateditem{Building Differential Homotopy Theory, Matsumoto, Japan}{Mar 2026}
  ``Geometric Quantization by Paths'' and ``The Boman Paradox''
  
  \dateditem{Shantou University, China}{Mar 2025}
  ``When Diffeology Meets Noncommutative Geometry''
  
  \dateditem{Duality Seminar, University of Warsaw}{Mar 2025}
  ``Why Diffeology''
  
  \dateditem{Fields Institute, Toronto, Canada}{Jul 2024}
  ``Symplectic Diffeology'' (Workshop on Hamiltonian Geometry and Quantization)
  
  \dateditem{Building Differential Homotopy Theory, Osaka, Japan}{Mar 2024}
  ``Symplectic Reduction in Diffeology with Singularities''
  
  \dateditem{Building Differential Homotopy Theory, Aizu, Japan}{Mar 2023}
  ``Symplectic Diffeology''
  
  \dateditem{Séminaire Interactions Fondamentales, CPT Luminy}{Dec 2022}
  ``Développements récents en diffeologie''
  
  \dateditem{Università degli Studi di Firenze, Italy}{Oct 2022}
  ``Every Symplectic Manifold is a Linear Coadjoint Orbit''
    
  %% Invited Lectures
  \subsection*{Invited Lectures (Earlier)} \null\
  
  \dateditem{Building Differential Homotopy Theory, Matsumoto, Japan}{Mar 2020}
  Series of lectures: ``Diffeology'' (invited course)

  %% Online Presence and Community Building
  \subsection*{Online Presence and Community Building} \null\
  
  \begin{itemize}
    \item \textbf{Website:} Founder, \url{Diffeology.net} – a central hub for the diffeology community.
    \item \textbf{Online Seminar:} Co-organizer, Monthly Global Seminar on Diffeology (2021–present).
    \item \textbf{Repository:} Creator of the \emph{Diffeology Archives} on GitHub – a definitive, version-controlled, machine-readable repository of foundational works in diffeology, designed for researchers and AI training.
  \end{itemize}
  
  %% Conference and Workshop Organization
  \subsection*{Conference and Workshop Organization} \null\
  
  \begin{itemize}
    \item International Conference Souriau 2019, Paris Diderot University (co-organizer)
    \item Workshop on Diffeology, Aix-en-Provence/Istanbul (2013–2016, series organizer)
    \item ``Nouvelles géométries pour la physique'', FRUMAM, Aix-en-Provence (2013)
    \item Souriau'90, International Conference in honor of J.-M. Souriau, Aix-en-Provence (2012)
    \item Summer School ``Symétries et application moment'', CIRM Marseille (1999)
    \item Le Séminaire des élèves, École Normale Supérieure de Lyon (1992–1998)
    \item Founder and Editor-in-Chief, \emph{Journal de Maths des Élèves de l'ENS-Lyon} (1994–1998)
  \end{itemize}
  
  %% Languages
  \subsection*{Languages}
  
  French (native), English (fluent), Hebrew (working knowledge)
  
  \vfill
  \noindent\textit{Last updated: May 2026. Version 1.0.}
  
\end{document}
